\chapter{Implementation of a Hybrid Simulator}
\label{chap:implementation}

TODO: reescrever este paragrafo
This chapter presents the implementation of this project's hybrid language, named Jaguar. It is organized into sections corresponding to each component described before: parser, interpreter, differential equation solver and visualizer. When relevant, formalizations are used to better detail the component being analyzed.

\section{System Architecture}
\label{sec:architecture}

TODO: passar coisas ditas em \autoref{sec:lince} para aqui.

\section{Parsing a Hybrid Language}
\label{sec:parser}

To parse our hybrid language, the Happy \todo{ref} parser generator was chosen. Happy takes an annotated BNF grammar specification and produces a Haskell module containing an LR parser for the grammar. Happy is the parsing library that GHC, the Haskell compiler, uses to parse its own syntax, and its pretty comprehensive functionalities proved to be more than enough for the project's needs.

%todo explicar gramatica (comentarios, precedencia de ops, chavetas, etc)

The full specification of the language's syntax, in BNF notation and its precedence rules, is in the appendice/user guide \todo{link}.


\section{A Hybrid Interpreter}
\label{sec:interpreter}

TODO: detalhes de implementacao (Hybrid, RunResult, etc)

\subsection{Iteration Limit}
TODO: zeno points
%todo



\subsection{Language Semantics: Formalization}
\label{sec:semantics}

%todo


\section{The Challenges of Undecidability}

Previously, we saw features ubiquitous among many hybrid simulators. However, dealing with computable real numbers results in some unique challenges, mainly when it comes to comparison. In this section, these challenges and their respective solutions are presented and used as motivation for the formalization of Jaguar's semantics in \autoref{sec:semantics}.


\subsection{Strict Comparison}
\label{sec:strict_comparison}

Take as an example the following program, which is meant to calculate the base 2 logarithm of a small number $x < 1$, rounded down ($\lfloor log_2 x \rfloor$):

\begin{lstlisting}[caption=Program to calculate the base 2 logarithm of x]
x := 0.0001;  //can be any number less than 1
y := 1;
ans := 0;
while x < y do { ans := ans - 1; y := y / 2 }
\end{lstlisting}

This algorithm is mathematically correct: it always takes a finite number of steps for any input $x$, and always reaches the correct answer. However, notice what happens if we try to run the program with the input $x=0.5$: the second iteration of the while loop will try to make the comparison $0.5 < 0.5$. Mathematically, it should evaluate to false, but as explained in \autoref{sec:theory}, comparing two equal computable real numbers will never terminate. As a result, our program will halt, even though it is mathematically sound!

The fundamental limitation here is the undecidability of comparison in the computable reals: whenever our hybrid simulator tries to fully evaluate a comparison, it may halt with no further feedback to the user. From a design standpoint, this places on the user the responsibility of ensuring that no comparison ever compares two equal numbers. Which, while in some cases is simply impractical, in others it is downright impossible.

Ideally, such comparisons would produce an error instead of halting. We can do that by evaluating the comparison only up to a certain precision. Formally, when evaluating $r<s$ to precision $n$, we first approximate $r$ and $s$ to intervals of rational numbers $a \leq r \leq b$ and $c \leq s \leq d$ such that $b-a,d-c \leq 2^{-n}$. If $b<c$, then we can be sure that $r<s$ and the original comparison is true. If $d \leq a$, then we are sure $r \geq s$ and the comparison is false. In all other situations, we can't resolve the comparison; the given precision isn't enough. All that can be inferred is that $|r-s| \leq 2^{-n+1}$. In these cases an error is raised, giving valuable feedback to the user. 

Analogous reasoning can be used for $>$, $\leq$ and $\geq$. However, operators for equality and inequality ($=$ and $\neq$) aren't included. The reason for this is simple: two equal numbers can't be guaranteed to be found equal, since this would require both numbers to be approximated to an interval of zero length (and while this isn't impossible - remember the restriction on intervals is $a \leq r \leq b$ such that $b-a \leq 2^{-n}$, so zero length is technically possible - it isn't guaranteed either, and depends entirely on the particular representation details of the two numbers). As a result, these hypothetical operators would always yield the same output: false in the case of $=$ and true in the case of $\neq$; or be unable to produce any output at all and raise an error.

Of course, if the comparison precision is too low, close comparisons ($|r-s| \leq 2^{-n+1}$ at most, depending on the approximations) can also result in an error; but these false positives are inevitable if we wish to prevent halting. It is up to the user to interpret the error, figure out the cause and decide how to solve it; either by raising the comparison precision, or by rewriting the program to avoid the faulty comparison, or even re-enabling full comparisons (by "raising the comparison precision to infinity") if they are confident the program will never halt.

The comparison precision $n$ could be specified for each comparison individually, but to avoid cluttering the syntax it was decided instead that all comparisons in a program would have the same precision, which can be supplied as an external configuration when running the program.


\subsection{Boolean Expressions}
\label{sec:bool_expr}

Since comparisons may not produce a result, boolean expressions using those comparisons have to deal with three possible outputs for a comparison: true, false and inconclusive. In the previous sections, it was implied that any inconclusive comparison would immediately raise an error, but this doesn't have to be the case. Take, as an example, the following program:

\begin{lstlisting}[caption=Hybrid program showcasing the usefulness of three-valued logic]
x := 1;
y := 1;
while x > 0.125 || y < 30 do { x := x / 2; y := y * 2; }
\end{lstlisting}

Although the comparison $x > 0.125$ is indeterminate at the beginning of the fourth iteration, the whole expression can be inferred to be true: since $y < 32$ is true, regardless of the value of the first term, the disjunction is true. In a sense, we can treat indeterminate values as simply another possible output from a comparison, and then work out the possible results of the disjunction. If both true and false results are possible, then the whole disjunction is indeterminate.

Essentially, what is being described is a three-valued logic as a way to capture the notion of indeterminacy. The topic is discussed in length in TODO\todo{ref}. Interestingly, this notion can also be conceptualized from a domain theory lens as the flat domain of booleans, containing three values: true, false and bottom. Such a formalization is useful because it allows us to extend what was said before about computable real numbers (\autoref{sec:domain}) to the semantics of Jaguar: evaluating boolean expressions will approach the actual result of the denoted condition, and expressions evaluated with greater comparison precision are always consistent with those evaluated with lower precision.\todo{proof?}

Introducing three-valued logic allows Jaguar to decide many expressions that were undecidable before, but is far from a full symbolic evaluation of the expression. For example, in an expression like $p \vee \neg p$, where $p$ is some comparison, should $p$ be inconclusive, the expression will be evaluated as inconclusive, even though logically it must always be true (the well-known tautology/contradiction problem of three-valued logic \todo{ref}).

Finally, notice that the three-valued logic is useless if the user re-enables full comparisons (up to "infinite precision"), since the comparison will simply halt instead of returning inconclusive.\todo{Unless all precisions are tried incrementally? Does this make sense?}

\subsection{Lax Comparison}

The solution described above for the undecidable comparison problem is as good as it gets: it detects faulty behavior and also allows the user to manage false positives. However, in some situations, the user may wish to disable errors entirely. Take as an example the following program, a simple while loop with an integer counter:

\begin{lstlisting}[caption=Simple "for" loop with an integer counter variable]
i := 0;
while i < 10 do { i := i + 1; }
\end{lstlisting}

The program is meant to exit the loop at 10 iterations, but the comparison $i < 10$ will raise an error when $i = 10$. A simple workaround is to replace the value 10 with a number we know $i$ will never be too close to and that still produces the same results when evaluating the comparison. In this case, $9.5$ works. But this puts the chore of finding such a number on the user, and the resulting code is much less readable (and much, much uglier).

But in this particular case, this replacement shouldn't be  necessary. Because we know $i$ is restricted to the integers, we can guarantee that an indeterminate comparison means that $i=10$ for any precision $n>1$, since the comparison is only indeterminate when $|i-10| \leq 2^{-n+1}$. Our restriction of $i$ gives us enough information to make these deductions, effectively bypassing the error.

For these situations, two new operators were created: determinably less than ($<!$) and determinably greater than ($>!$). Unlike the strict comparison operators (the ones introduced in \autoref{sec:strict_comparison} that can return true, false or indeterminate), these new lax comparators only return true or false: true if the comparison is determinably true, and false in all other cases. They are lax in the sense that they always produce an answer, even when there isn't enough information.

This opens the door to all sorts of misbehavior in the program, of course: $9.99 <!\ 10$ will evaluate to false if the comparison precision is too low. Indeed, the output of this operator may be incorrect for all comparisons $r <!\ s$ where $r < s \leq r + 2^{-n+1}$. However, outside this range, it is guaranteed to be correct. So long as this critical range is avoided, the lax comparators can be used safely. If both operands are integers, this works out to $n>1$, as stated above. For other cases, it may vary. It is the responsibility of the user to make sure the variables being compared are appropriately constrained and choose the appropriate comparison precision for their use case.



\subsection{Runtime Errors}

As seen in the previous sections, our language can throw errors. \todo{semantics of errors?}
%todo: operational semantics of error? (the composition of two programs is an error if either of the two subprograms is an error)

To handle errors without raising a native Haskell exception, the output of a program can be either a valid state or an error state, containing an error message. The following is a list of all situations that can raise an error:
\begin{itemize}
  \item An inconclusive boolean expression, caused by a strict comparison (Comparison precision exhausted)
  \item Too many iterations of all while loops, summed together (Iteration limit exceeded)
  \item A differential statement with a non-positive or inconclusive time step
\end{itemize}


\subsection{Hybrid Undecidability}
\label{und}

Another unique problem to the crossing of hybrid systems and arbitrary-precision real numbers is the issue of undecidability on the concatenation point of two hybrid programs. Take the following program:

\begin{lstlisting}
x:=0;
wait 1;
x:=1;
x' = 1 for 1;
\end{lstlisting}

The output of this program should be a function $x(t)$ such that:

\begin{equation}
  \begin{cases} 
    x(t) = 0 & \text{if } 0 \leq t < 1 \\
    x(t) = t & \text{if } 1 \leq t \leq 2 \\
  \end{cases}
\end{equation}

Although this function is mathematically well-defined, in order to evaluate it to any particular value of $t$ we must perform the comparison $t < 1$ to decide between the two branches of the function. However, because comparison is undecidable on the computable reals, evaluating the function may be a halting computation for $t=1$, and may take arbitrarily long when evaluated for values close to $1$.

To solve this problem, as before, we must use a parametrized comparison, taking in the desired precision of our query. Note that this precision doesn't necessarily need to match the comparison precision used for in-program comparisons. The query precision is only used to choose between two (or more) consecutive hybrid states of the program, leaving the actual execution of the program (variable assignments, choosing between if-else branches and while-loop iterations) completely unaffected. Indeed, the query precision can even be chosen after the whole program has already run, and can be chosen independently for every query.

Taking that into account, to give the user the most flexibility and control, it was decided that a query would be a curried function $\mathbb{R} \rightarrow \mathbb{N} \rightarrow [\mathbb{R}^n]$ that takes the time $t$ and the desired query precision $n$ and outputs all possible hybrid states within range of that precision. The biggest drawback of this solution is that it may lead to some function branches being evaluated outside their original domain (in this particular example, evaluating the query $t=0.999$ with a query precision $n < 10$ will result in two possible outputs: $0$, the correct output; and $0.999$, which is clearly impossible to obtain from the mathematical definition, but must be considered nonetheless).

\section{Differential Equation Solver}
\label{sec:solver}

%introduction (IVP, numeric vs symbolic)
As explained above, the continuous behavior of our hybrid programs is modeled using differential equations. This section exposes how to use a differential equation to evolve a set of initial conditions - the well-known initial value problem.

%adapted from
%https://www.maia.ub.es/~angel/taylor/taylor-1.4/taylor.pdf
The general form of the initial value problem (IVP \todo{indice de abreviaturas?}) can be formalized as follows: given the initial value of a smooth unknown function $x: [a, b] → R^m$ at $a$ such that

\begin{equation}
  \begin{cases} 
    x(a) = x_0 \\
    x'(t) = f(t, x(t))
  \end{cases}
\end{equation}

with $f: [a, b]×C ⊂ R×R^m → R^m$, $m≥1$ and $C$ a closed set, find the value of $x(t)$ for any $t \in [a,b]$.

As discussed previously, our solver should work numerically to take advantage of the operations of computable real numbers. In general, numerical methods for solving the initial value problem involve calculating or estimating the first (and/or higher) order derivative(s) of $x$ at $a$ and then estimating $x(a+h)$, for some step size $h$, with a formula that uses the obtained derivative(s). This logic comes from [... taylor theorem...]. Different methods are then a balancing act between using many derivatives or small step sizes, which are more accurate but more expensive computations, and few derivatives or big step sizes, which are less accurate but much faster to compute, all while preventing compounding errors from all of the estimating. [... exemplos e citacoes]

This puts us in an awkward position: since using any sort of estimation compromises the arbitrary precision we are aiming for, most of these methods are useless to us. The arbitrary precision of computable real numbers requires two things: that the precision of the solution can be increased on demand; and that the solution has a strict error bound. And, ideally, it should be done with the least numerical operations possible, since operations on computable reals are much slower than their floating-point counterparts.

These restrictions naturally lead us to the Taylor method. Formally, the Taylor method can be described as [...TODO...] By using the exact values of the derivatives, the only error introduced is from the truncation of the potentially infinite Taylor series. Acquiring a more accurate approximation is as simple as increasing the order $p$, which simply means more terms are used, thus reducing the truncation error. Calculating the error bound of the approximation is discussed below.

To implement the Taylor method, we therefore need to calculate the derivatives of $x$. Although obtaining the first-order derivative is a simple calculation ($x'(a) = f(a, x(a))$), TODO

%taylor method, that segways into
%automatic differentiation (implementation in haskell???)
%formally define problem: find numerical solutions to systems of ODEs (restricted to some operations) with strict error bounds
%compare possible methods (interval taylor, polynomial)

%ref: como transformar ODE em ODE polinomial: https://educ.jmu.edu/~sochacjs/PolynomialODEs-Sochacki.pdf

%matematica do solver - ou seja, pegar no Teorema 3 (TM) e derivar as minhas formulas de accuracy
%falar sobre solucoes polinomiais vs series de taylor infinitas e k-skip


\section{Visualizer}
\label{sec:visualizer}

\subsection{Discontinuity Tracking}

It was already explained in \autoref{und} how discontinuities (which may arise at the point of concatenation of two hybrid programs) are handled locally, when making individual queries. However, to facilitate a big picture analysis of a program (useful when plotting the system trajectory), Jaguar also includes tools for tracking discontinuity points throughout the whole trajectory. This is as simple as keeping track of every point at which an assignment is made\footnote{It is assumed that differential statements do not generate discontinuities in the mathematical sense, since analyzing the differential equation for discontinuities goes beyond the scope of the project}. Not every assignment generates a discontinuity, since it may not change the state, but checking if the state changed requires comparisons (which, as seen before, leads to problems of comparison precision and undecidability), so all assignments are tracked. The resulting list of (potential) discontinuities is returned alongside the evolution function as the output of a program.

Each discontinuity is located between two differential trajectories. In order to allow matching each discontinuity to its neighboring hybrid branches without directly comparing the time of each (which, again, is undecidable), each discontinuity and differential trajectory also include order information in the form of a step counter. Essentially, alongside the state of the system, Jaguar keeps a counter that starts at zero and increases with both differential statements and assignments. This way, each differential trajectory is associated with a unique step number.

\todo{example}
%todo: example program, image of plot (with the step counter added on top) and list of discontinuities

A discontinuity is described by all that information put together: a real number corresponding to the time of the discontinuity; the hybrid state and step number "before" the discontinuity (as in, the state the system would have at the moment of discontinuity if the program ended without the assignment statement that caused the discontinuity); and the hybrid state and step number at the moment of the discontinuity.





%todo: fazer features novas (cmd line query mode, interactive plot)
%   antes disto tudo, explicar CompReal, approx, lower e limit
%   discs e Unds
%   desvantagem: a biblioteca que eu estou a usar usa Doubles, logo muito zoom pode causar problemas (todo: resolver com novo plot interativo?)
